\section{精确陈述}

我们首先以 Lean 开发实际证明的形式陈述抽象定理。这防止几何直觉在不知不觉中加强结论。

固定流形标签集合、有理向量空间分次族 $G(M,q)$、标签 $X$ 与 $S^4$，以及谓词
\(\operatorname{HS}(M)\) 与 \(\operatorname{Diff}(M,N)\)。

\begin{theorem}[抽象非标准性定理]\label{thm:A}
设有有理向量空间 $W_0,W_1,W_2,W_3$、元素 $x_0\in W_0$、线性映射
\[
W_0\xrightarrow{q_{01}}W_1\xrightarrow{q_{12}}W_2,
\qquad
\ell_i:W_i\longrightarrow\Q\quad (i=0,1,2),
\]
以及线性同构
\[
T:W_2\xrightarrow{\cong}W_3,
\qquad
F:W_3\xrightarrow{\cong}G(X,494),
\]
使得
\begin{align}
\ell_0(x_0)&=2624,\label{eq:value}\\
\ell_1q_{01}&=\ell_0,\label{eq:q01}\\
\ell_2q_{12}&=\ell_1.\label{eq:q12}
\end{align}
再假设 $G(S^4,494)$ 的每个元素为零，且每个微分同胚 $X\to S^4$ 诱导线性同构
\(G(X,494)\cong G(S^4,494)\)。则 $X$ 不与 $S^4$ 微分同胚。
\end{theorem}

\begin{proof}
令 $v=F(T(q_{12}(q_{01}(x_0))))$。若 $v=0$，$F$ 与 $T$ 的单射性蕴含 $q_{12}(q_{01}(x_0))=0$。应用 $\ell_2$ 并用 \eqref{eq:q01}--\eqref{eq:q12} 得 $2624=0$，矛盾。故 $v\ne0$。若 $X$ 与 $S^4$ 微分同胚，诱导同构将 $v$ 映到 $G(S^4,494)$ 的元素，故为零，与单射性矛盾。
\end{proof}

与定理~\ref{thm:A} 对应的 Lean 定理是从输入结构 \texttt{ExternalGeometry} 与
\texttt{CSExternalGeometry} 出发的条件蕴含；Lean 中不构造其实例。
候选特定输入见下文，并在第~\ref{sec:johnson-presentation}--\ref{sec:final-identifications} 节证明。

\begin{theorem}[嵌入候选表示，P0]\label{hyp:P0}
显式 Johnson 生成元替换是 $\Sigma_A^0$ 的带标号带框柄表示，含两个带框乘积消去、嵌入 detector 球及所述缆附着。
\end{theorem}

\begin{proof}[定理~\ref{hyp:P0} 的证明]
见第~\ref{sec:johnson-presentation} 与~\ref{sec:p0-conclusion} 节，特别是定理~\ref{thm:P0discharge}。
\end{proof}

\begin{theorem}[系数比较，C]\label{hyp:P1}
完备 MWW 系数量子 trace 允许到 BPW/BHPW 权重一 endpoint 表示的分次保持、对称幺半自然映射，对远离 $B$ 支持的边界 cobordism 亦自然。在该映射下 point-push 作用为公开 Burau 表示，divided cubic 行通过每个 beta 与 psi 关系下降。其在所选类上的值为 $2624$。
该类位于 Johnson 替换领圈的 quantum degree $494$。
\end{theorem}

\begin{proof}[定理~\ref{hyp:P1} 的证明]
见第~\ref{sec:c-comparison} 节，特别是定理~\ref{thm:Cdischarge}。
\end{proof}

\begin{theorem}[相对三柄闭包，S]\label{hyp:P2}
三个 $3$-柄可沿与 $B$ 不相交的完备球面系统附着。对每个球面，detector 消去无点 MWW 关系并将一次带点映射与源项等同。故 divided detector 通过完备三柄 coequalizer 下降。
\end{theorem}

\begin{proof}[定理~\ref{hyp:P2} 的证明]
见第~\ref{sec:s-closure} 节，特别是定理~\ref{thm:Sdischarge}。
\end{proof}

\begin{theorem}[替换图景的四柄闭包，P3]\label{hyp:P3}
在假设~\ref{hyp:P2} 的反向 $1$-柄 $3$-柄之后，剩余边界为 $S^3$。沿该 $S^3$ 附着 $4$-ball 得闭 $4$-流形 $X_J$。MWW 命题~3.4 将该 $W_3$ 图景的空 link lasagna 模与 $\mathcal S^2_0(X_J)$ 等同。MWW 推论~3.5 给出 $\mathcal S^2_0(S^4)\cong\mathbb Z$ 集中在 bidegree $(0,0)$，故 quantum-$494$ 分项消失。微分同胚诱导分次保持等价。$\det A=\det(A-I)=1$ 是标准 Cappell--Shaneson 构造的同伦球面判据。识别 $X_J\cong\Sigma_A^0$ 是假设~\ref{hyp:P0} 的构造 Johnson CS 柄图，而非这些行列式。
\end{theorem}

\begin{proof}[定理~\ref{hyp:P3} 的证明]
反向 $1$-柄图景在三个 $1$--$3$ 消去后边界为 $S^3$；附着 PL $4$-ball 得闭流形 $X_J$。空 link Khovanov 在 quantum degrees $0$ 与 $494$ 的秩、附着 $4$-ball 的 Euler 特征，以及分次和 $-44+227+315-4=494$ 均直接计算。同构 $\mathcal S^2_0(B^4)\cong\mathcal S^2_0(S^4)$ 与 $\mathcal S^2_0(S^4)\cong\mathbb Z$ 集中在 bidegree $(0,0)$ 见 \cite[Proposition~3.4 and Corollary~3.5]{ManolescuWalkerWedrich2023}。识别 $X_J\cong\Sigma_A^0$ 是定理~\ref{thm:P0discharge}。独立重放命令列于附录~\ref{sec:reproducibility-extended}。
\end{proof}

\begin{theorem}[条件 trace-73 定理]\label{thm:joined}
Johnson CS 柄图 $X_J$ 与 $\Sigma_A^0$ 等同。Iwaki 命题~2.1 因此使 $\Sigma_A^0$ 为同伦球面。若假设~\ref{hyp:P0}--\ref{hyp:P3} 的几何数据组装进 Lean 接口 \texttt{ExternalGeometry}，则
\[
\Sigma_A^0\not\cong_{\mathrm{diff}}S^4.
\]
该接口未在 Lean 中构造，故不断言反例。
\end{theorem}

\begin{proof}[定理~\ref{thm:joined} 的证明]
\texttt{ExternalGeometry} 有实例时，定理~\ref{thm:A} 给出光滑阻碍。识别 $X_J\cong\Sigma_A^0$ 与定理~\ref{thm:P0discharge}、\ref{thm:Cdischarge}、~\ref{thm:Sdischarge} 提供假设。
\end{proof}

\input{sec-published-results-zh}
